\section{Introduction}
\label{sec:intro}

Fabric defect detection is a representative industrial visual inspection task
that demands high recall, low latency, and low computational cost
simultaneously. In production settings, missing a defect is costly, yet
inspection must keep pace with high-throughput manufacturing on commodity
hardware. A central difficulty is that many fabric defects---nep, broken or
hanging warp, weft shrinkage, and weak-texture anomalies---are extremely small
or visually faint, and are easily missed when detectors operate at the low
input resolutions favored for efficiency.

Most prior work treats such misses as a model-capacity problem, addressing them
through network architecture, loss design, knowledge distillation, multi-scale
training, or sliding-window inference. These approaches implicitly assume that a
missed defect could have been detected with a better low-resolution model. They
do not distinguish between a \emph{recoverable} miss---where the evidence is
present but the model is insufficiently sensitive---and a miss caused by
\emph{insufficient input information}, where the defect occupies too few pixels
for any low-resolution model to recover.

We make this distinction concrete through a cross-resolution diagnosis. Running
the same detector at low ($640$) and high ($1280$) resolution partitions missed
defects into three regimes: threshold-suppressed misses recoverable by
calibration; visibility-limited misses recoverable only by genuine high
resolution; and ambiguous misses that persist even at high resolution. We then
show empirically that a range of cheap interventions---per-class threshold
calibration, miss-rate loss reweighting~\cite{focalloss,eql}, a small-object detection head, feature
distillation, and sparse patch reinspection---fail to recover the
visibility-limited misses. Only full-image high-resolution inference recovers
them, but it quadruples cost and, on easy data, can even reduce accuracy.

This motivates treating the problem not as one of improving low-resolution
representation, but as one of \emph{deciding when to pay the high-resolution
cost}. We propose Visibility-Boundary-Aware Detection (VBAD). VBAD constructs a
visibility label automatically from cross-resolution failure agreement
(low-resolution miss followed by high-resolution recovery), trains a lightweight
visibility head on a frozen detector to predict it, and selectively triggers
full-image high-resolution reinspection only on images judged
visibility-limited. It additionally surfaces the ambiguous, unrecoverable
fraction as a reliability boundary for human review.

Our contributions are:
\begin{enumerate}[leftmargin=1.5em]
\item A cross-resolution failure diagnosis for industrial fabric defect
detection, showing that missed defects arise from threshold suppression,
visibility limitation, and ambiguous low-reliability cases, and that cheap
low-resolution interventions empirically fail to recover the visibility-limited
ones.
\item A visibility-boundary-aware detection framework that learns from
cross-resolution failure supervision to selectively trigger full-image
high-resolution reinspection, improving the recall--cost trade-off over both
uniform low- and high-resolution inference.
\item Validation on a public Science Data Bank fabric defect dataset, including
an ambiguity boundary analysis that quantifies how many missed defects are
inherently unrecoverable and should be escalated for manual review.
\end{enumerate}
